\chapter{Conclusions}
\label{chap:conclusioni}

\section{Achieved results}

\autoref{chap:introduzione} sets the objective of designing, implementing, and validating RAM-USB as a
multi-user remote backup service compliant with the requirements of the SRS \cite{ramusb-srs}, relying on the
zero-trust and defense-in-depth principles to defend the system from external attacks.

The ten user requirements listed in \autoref{chap:requisiti} are, with only two exceptions, fully
covered by
system requirements implemented and traced all the way to the code and the tests. A user can register by providing
only an email, a password, and an SSH public key (RU-01), authenticate in subsequent sessions (RU-02),
upload and retrieve client-side encrypted backups from anywhere (RU-03, RU-07), do so only after having
authenticated (RU-04), with the guarantee that the content of the backups is unreadable even by
an administrator (RU-05), in a space isolated from that of other users (RU-08).
An administrator can observe the health and performance of the system without ever touching individual
users' data (RU-10).

RU-06 (Have the guarantee that no one, administrators included, can read one's own
access credentials) is not fully supported because of the theoretical limitations imposed by the choice of an email and password
authentication system, and therefore holds only for the password and only once it is already
persisted in the database in hashed form. The only way to prevent an administrator from intercepting
credentials in transit, in fact, is not to send them to the server at all, for example by using a
zero-knowledge protocol, which is not the subject of this thesis.

RU-09, the guarantee that no one can modify or delete a backup, was
declared out of scope from the earliest versions of the SRS, because of the project's tight schedule, and is
taken up again in the next section.

The five use cases that run through the system (registration, login, backup, restore, and
consulting the metrics) were executed end-to-end on the real Proxmox deployment, with the
same results reported in \autoref{chap:testing}.

\section{Known limitations}

The SRS explicitly accepts four risks as the cost of releasing on schedule \cite{ramusb-srs}.
Normally these risks would not be acceptable in a real backup system, and their resolution is
planned as future work.

RISK-01 is RU-09 itself. It remains a desirable improvement, not yet designed, implemented, or tested.

RISK-02 concerns the master key and the pepper: they reside in an environment variable, with no backup procedure
(DV-F-18) or rotation procedure (DV-F-19). Their loss
would entail the irreversible loss of access to all encrypted credentials.

RISK-03 concerns the client: it is currently not containerized with Docker. The service's Go code exists and has
been tested, but containerization requires further reasoning and design decisions.

RISK-04 is the circular dependency of Network-Manager's DNS resolution on its own MagicDNS: if
Network-Manager published a non-working ACL policy, it could also lose the ability to
resolve the Headscale hostname it would need to correct it. In that case, it is up to an
administrator to resolve the problem manually.

\section{Future work}

\autoref{sec:entropia} shows that, under uniform choice, the constraint of at least three character
categories is worth less than a tenth of a bit of entropy, while one additional character is worth 6.57 bits: under
that assumption the composition rule is redundant, because a truly random password
already uses several categories by itself. Its real purpose, however, is not to increase entropy but to constrain a human's
choice, and in that respect its value cannot be measured without the real distribution of the passwords chosen
by users.
This design choice therefore currently remains open, pending more in-depth analysis.

Besides DV-F-18 and DV-F-19, already discussed as RISK-02, two items remain to be designed: the per-user
storage space limit and the automatic replication of backups on CephFS, both
deferred in the SRS as desirable improvements rather than project constraints.

Other future work, not yet designed or included in the SRS, comprises automatic load
balancing, protection of Entry-Hub's public endpoint from availability attacks,
perhaps by imposing rate limiting on that endpoint, and improved system scalability, perhaps
by unlocking geo-distribution of the server.
